% !TeX root = ncaa-manuscript-en.tex
\section{Discussion}\label{sec:discussion}

The discussion follows the three problems in the Introduction: constructing
and checking source-bound constraints, preserving their meaning in learning
sentences, and using them in action selection. The distinction between
constraint compliance and preferences among admissible candidates applies
throughout; a high preferred-selection score concerns the supplied candidate
set and objective.

\subsection{Constraint Construction and Specification Checkability}\label{sec:discussion-checkability}

Reinforcement learning can improve constraint compliance and action selection, but
the rules and evaluation criteria used for learning require separate validation:
a misspecified rule or reward can yield a policy that optimizes the supplied
signal while departing from the intended requirements. In a traffic-merging
simulation, Pan et al. report that an agent optimizing a proxy reward based on
average vehicle speed learned to stop its vehicle instead of merging, increasing
the proxy reward while also increasing mean commute time~\cite{Pan2022RewardMisspecification}.
When several rules and modeling assumptions apply, their joint consistency and
whether their combination permits necessary safety responses also require
examination. This need to check the learning specification applies to both
supervised and reinforcement learning.

The EBLC path links each constraint field to its source, making mismatches
traceable and available for review. A constraint can draw on an observed
clearance time, a supplied waiting duration, and an assumption that the zone
remains clear. A sentence alone may not identify which of these inputs is
responsible for an error. EBLC separates them into source-linked fields that
can be checked individually.

The composition diagnostic adds a distinct question: can conditions from
separate source records hold together? Its three-way braking conflicts show
that individual consistency, and even consistency of every pair, can leave a
joint incompatibility undetected (Section~\ref{sec:results-composition}).
Identifying a subset-minimal source group makes the requirements needing
review explicit. The source roles separate documented policy or guidance from authored
observations, candidate requests, objectives, and modeling assumptions;
their applicability is supplied rather than inferred. A SAT combination establishes compatibility under the specified
context, while an UNSAT combination calls for reviewing requirements or
assumptions rather than silently discarding one. Hold and non-crossing controls show why permission and progress obligations
must be distinguished: waiting may satisfy restrictions even when a proposed
movement is blocked. Comparing the faithful and omitted-exception Rule 129
encodings further identifies a lost permission, rather than a conflict in the
original regulation; detecting an omitted exception therefore requires a
reference encoding that retains it.

P1 is the natural-language guard path of our preceding work~\cite{r4}, and P2
is the EBLC path introduced here. The comparison therefore concerns two
development paths for the same constraints. Both achieve high primary-test
preferred-selection scores, and neither dominates every condition
(Sections~\ref{sec:results-six-tasks}--\ref{sec:results-robustness}). P2 is
ahead in position-hold continuation with XY coordinates under the 2.0~s rule
(97.92\% versus 91.15\% preferred-selection score; 4 versus 16 violations of 192) and on
pedestrian-entry paraphrases (96.94\% versus 90.56\%), whereas P1 is ahead on
static paraphrases and the combined maneuver set. Because the paths also differ
in wording, these differences are not attributed to specification checking. The EBLC path
thus retains the behavioral benefit of constraint training and adds the ability
to check constraint construction and correspondence.

Assurance in this process can be considered from three perspectives; the
first two group the five checks in Section~\ref{sec:method-checks}:
\begin{enumerate}
\item \textbf{Specification structure and logical correspondence.} Does the
specification express the intended restriction and release within the supported
model, rules, and queries? Structural validation, satisfiability checks, and
logic--rule comparisons assess this question.
\item \textbf{Evidence and sentence correspondence.} Do specification fields
match their source records, and does parsing the generated sentence with the
fixed grammar recover those fields? Evidence and sentence correspondence
checks assess these relationships.
\item \textbf{Model behavior.} Does the model select an admissible candidate,
and does it select the preferred candidate under the stated objective?
Constraint violation rate and preferred-selection score measure these two
questions separately.
\end{enumerate}

Fault injection shows complementary roles for the tested checks:
specification-level checks detect faults in the specification and logical
formula, whereas only sentence correspondence detects faults introduced in the
generated CNL (Table~\ref{tab:results-faults}). Evidence and sentence
correspondence checks, together with boundary probes for logic--rule comparison,
were added for this evaluation after the reported training runs.

Establishing the validity of a source record requires independent observations
or rule evidence. All 48 controls in which the record and specification field
are changed to the same wrong value remain internally consistent and escape
every check. A possible division of work in practice would concentrate human
review on the suitability of observations, rules, and assumptions, while
automating structural, logical, and sentence-correspondence checks.

\subsection{From Checked Meaning to Learned Interpretation}\label{sec:discussion-interpretation}

Specification checking and behavioral evaluation answer different questions:
what the constraint says and how the model uses it. Recovering
specification fields from a generated sentence does not establish how the
model interprets that sentence. The results illustrate this distinction:
\begin{itemize}
\item \textbf{Paraphrasing.} Changing wording while retaining the intended
rule reduces static-task P2 preferred-selection score from 100\% to 80.32\%; P1 reaches
93.29\% under the same condition (Table~\ref{tab:results-robustness}).
\item \textbf{Logical text.} Primary preferred-selection score is 99.17\%, but on unseen waits
it is 81.25\% despite only 0.83\% violations
(Section~\ref{sec:results-robustness}).
\item \textbf{Combined difficulty.} Combining negation, clause order,
equality boundaries, and units reduces P2 rule-pair accuracy from 97.92\%
to 82.64\% (Section~\ref{sec:results-robustness}).
\item \textbf{Image dependence.} Although candidate values are supplied in
text, a blank image reduces P2 preferred-selection score from 98.61\% to 76.39\%
(Section~\ref{sec:results-visual}). These interventions measure input dependence; visual grounding and event
interpretation require separate assessment.
\end{itemize}

Multiple fixed expressions designed to recover the same specification fields
offer a starting point for combining checkability with expression diversity.
A fixed grammar enables sentence-correspondence checks, but a restricted set
of training formulations may encourage dependence on particular wording.
Whether training jointly with these expressions reduces wording sensitivity
remains to be evaluated. Such checks apply within the defined grammar; they
do not determine semantic equivalence for arbitrary natural-language sentences.

\subsection{Learning Constraint-Conditioned Candidate Selection}\label{sec:discussion-training}

Under the evaluated supervised-learning conditions, with initialization and
evaluation inputs held fixed, training with constraint
sentences improves rule-dependent candidate selection. In the common-CNL
evaluation (Table~\ref{tab:results-summary}), CE adapters trained with CNL
reach 96.94\% pedestrian-entry rule-pair accuracy, compared with 1.94\% for
CoC-only-trained adapters. This contrast concerns selecting both reference
labels when the scene and candidates are identical but the rules differ:
training with constraint text improves the use of a rule available at
evaluation relative to adapters fine-tuned in a rule-free format.
Position-hold results show the same pattern.

Complementary evidence points in the same direction, although these
evaluations do not isolate the training effect. Without task-specific training,
the 8B model violates the supplied pedestrian-entry rule in 68.75\% of
predictions (Table~\ref{tab:results-pretrained}). The preliminary Alpamayo-R1
trajectory diagnostics likewise exceed supplied speed limits in 15 of 18
outputs, with the same violation counts whether the rule is supplied through
the user prompt or appended to CoC
(Section~\ref{sec:introduction}; Supplementary Section S1). Together, these
findings indicate that the presence of rule text at evaluation is insufficient
for strong rule-dependent selection in the tested models, whereas training
with constraint text in the evaluation format substantially improves that
selection in the controlled study.
This parallels instruction tuning, in which training with task instructions
improves inference-time task performance~\cite{Wei2022}; here, the instruction
is a case-specific constraint with numerical parameters, and its effect is
tested by changing only the rule.

Providing constraints through the input and through the loss has different
effects. CE supervises the reference label of the preferred candidate, while the semantic term in
Eq.~\ref{eq:setup-semantic-loss} increases the conditional probability mass
assigned to admissible labels. Without rule text, the same scene and candidates
can have different reference labels under two rules, leaving the inputs
unable to distinguish the cases. Adding semantic loss nearly eliminates
violations, but preferred-selection score remains near the rule-free ceilings of 50\% and
75\% (Table~\ref{tab:results-semantic}). This is consistent with more
conservative choices admissible under both rules. Loss-based constraint
learning encourages satisfaction~\cite{Xu2018,Giunchiglia2022}, whereas the
input rule identifies which constraint applies. Adding the tested semantic
term when rule text is already supplied gives similar results overall.

Reinforcement learning can also train a rule-conditioned policy when the
applicable rule is available in its input. The present CE and semantic-loss
comparisons establish the effect of constraint text within supervised
fine-tuning; they do not establish superiority over reinforcement learning.
Checking the supplied constraints remains a separate concern whichever
optimizer is used, as discussed in Section~\ref{sec:discussion-checkability}.

Our learning objective is to select actions conditional on the rule supplied
for each example. Constraints such as waiting durations or distance limits
may differ even for the same scene because of local regulations, temporary
signs, or operating instructions. The intended behavior is to receive the
applicable constraint and adjust the choice accordingly. Deliberative Alignment
trains models to recall reusable safety policies and apply them in context;
policy text is used for data generation and filtering, then removed from
supervised fine-tuning inputs~\cite{Guan2024}. Here, example-specific constraints and their numerical
parameters remain in both training and evaluation inputs, and the evaluation
tests whether choices respond to changes in those constraints.

\subsection{Limitations and Future Work}\label{sec:discussion-limitations}

\noindent\textbf{Tasks and data.}
The learning experiments select from finite candidate sets in synthetic tasks.
Candidate values are supplied in text, and static and maneuver diagrams are
reused across splits. The results provide controlled evidence about constraint
training; extending them to real-road visual understanding, trajectory
generation, or closed-loop safety requires separate evaluation. The Alpamayo-R1
diagnostics motivate the problem and do not constitute results from training
Alpamayo.

\noindent\textbf{Constraint compliance near boundaries.}
Selection improves on the new position-hold test, but the earlier XY regression
diagnostic retains 106, 79, and 77 violations out of 192 predictions for P0,
P1, and P2 under the 2~s rule after training
(Section~\ref{sec:results-continuation}; Supplementary Section S11). This
diagnostic includes prohibited candidates only 1~mm beyond the radius, compared
with a minimum overshoot of 22~mm in the new test. The results therefore do
not establish stable constraint compliance across conditions involving small overshoots.

\noindent\textbf{Comparison design.}
Fine-tuning uses one 2B model family and three optimization seeds. Continuation
comparisons include differences in prior training history; the fresh-adapter
study separately controls initialization. CoC-only fine-tuning uses a rule-free
input format, whereas common-CNL evaluation adds text absent during that
fine-tuning. The comparison therefore includes adaptation to the constraint
input format; it does not isolate how much of the gap arises from learning
rule-conditioned selection versus reduced use of additional text after
rule-free fine-tuning. The tested semantic loss acts only
on conditional A--D label probabilities at the first supervised answer token,
uses a weight fixed at $\lambda=1$ before execution, and leaves alternative
formulations to be compared. It provides a training penalty rather than a
generation-time restriction or a vehicle-safety guarantee. The pretrained 8B
evaluation includes task-format adaptation.
No equivalence or noninferiority test between P1 and P2 was conducted. Their
differences reflect wording and construction paths and do not isolate the
causal effect of specification checking.

\noindent\textbf{Checking.}
Because evidence correspondence, sentence correspondence, and boundary probes
for logic--rule comparison were added after training, their effects on training-data quality and model behavior remain to
be measured. Detection results concern supported grammars and
injected fault types; the prevalence of naturally occurring errors must be
estimated from actual constraint-authoring data. Evidence binding separates
observations, rules, and assumptions by design, but the injection study did
not test incorrect scene observations or unsuitable modeling assumptions; its
faults concern rule values, logical restrictions, units, and generated CNL.
The composition diagnostic uses a fixed, manually constructed grid grounded
in selected document clauses. Authored states, simplified braking and signal
models, candidate obligations, and one restricted exception define its scope.
It does not evaluate complete regulations, emergency-trigger eligibility,
jerk-limited braking, learned reward models, or operational safety. Its
natural-language statements explain the formalization, rather than testing
automatic legal interpretation or composition of CNL training examples. The accident examples
motivate those broader checks; they are not evidence that this implementation
would have prevented the accidents.

Future work should:
\begin{enumerate}
\item Evaluate how observations and rule applicability are assessed and bound
in real scenes.
\item Apply constraint-bearing training data to trajectory-generating VLA
models and measure trajectory-level constraint compliance, including near boundaries.
\item Evaluate deployment-time generation of constraint sentences from scene
observations and applicable rules for provision to the model.
\item Test whether training with multiple expressions that recover the same
specification fields reduces wording sensitivity.
\item Repeat the controlled fresh-adapter comparison with VLMs of other
families and sizes.
\item Measure human review time and error rates for constraint authoring in
an actual workflow.
\end{enumerate}
